% Part: normal-modal-logic
% Chapter: filtrations
% Section: finite

\documentclass[../../../include/open-logic-section]{subfiles}

\begin{document}

\olfileid{nml}{fil}{fin}

\olsection{تصفیه‌ها متناهی‌اند}

تصفیه‌ها را برای هر مجموعهٔ~$\Gamma$ که تحت زیر-!!{formula}ها بسته
باشد تعریف کرده‌ایم. هیچ‌چیز در خود تعریف تضمین نمی‌کند که تصفیه‌ها
متناهی باشند. در واقع، وقتی~$\Gamma$ نامتناهی است (برای نمونه، مجموعهٔ
همهٔ !!{formula}هاست)، تصفیه نیز می‌تواند نامتناهی باشد. بااین‌حال، اگر
$\Gamma$ متناهی باشد (برای نمونه، هنگامی که مجموعهٔ زیر-!!{formula}های
یک !!{formula} مفروض~$!A$ است)، هر تصفیه‌ای از خلال~$\Gamma$ نیز
متناهی است.

\begin{prop}\ollabel{prop:filt-are-finite}
  اگر~$\Gamma$ متناهی باشد، آنگاه هر تصفیهٔ~$\mModel{M^*}$ از مدل
  $\mModel{M}$ از خلال~$\Gamma$ نیز متناهی است.
\end{prop}

\begin{proof}
  اندازهٔ~$W^*$ شمار رده‌های گوناگون~$[w]$ تحت رابطهٔ هم‌ارزی
  $\equiv$ است. هر دو جهان~$u$ و~$v$ در چنین رده‌ای---یعنی هر~$u$ و
  $v$ که~$u \equiv v$---دربارهٔ همهٔ !!{formula}های~$!A$ در~$\Gamma$
  توافق دارند: برای~$!A \in \Gamma$، یا~$!A$ هم در~$u$ و هم در~$v$
  صادق است، یا در هیچ‌یک. پس هر ردهٔ~$[w]$ با زیرمجموعه‌ای از~$\Gamma$
  متناظر است؛ یعنی مجموعهٔ همهٔ~$!A \in \Gamma$ که در جهان‌های~$[w]$
  صادق‌اند. هیچ دو ردهٔ متمایز~$[u]$ و~$[v]$ با یک زیرمجموعهٔ واحد از
  $\Gamma$ متناظر نیستند. زیرا اگر مجموعهٔ !!{formula}های صادق در~$u$
  و مجموعهٔ !!{formula}های صادق در~$v$ یکسان باشند، آنگاه~$u$ و~$v$
  دربارهٔ همهٔ فرمول‌های~$\Gamma$ توافق دارند؛ یعنی~$u \equiv v$.
  اما در این صورت~$[u] = [v]$. بنابراین تابعی !!a{injective} از~$W^*$
  به~$\Pow{\Gamma}$ وجود دارد، و ازاین‌رو
  $\card{W^*} \le \card{\Pow{\Gamma}}$. پس اگر~$\Gamma$ شامل~$n$
  جمله باشد، کاردینالیتی~$W^*$ بیش از~$2^n$ نیست.
\end{proof}

\end{document}
